\chapter{O $\lambda$-cálculo}

\section{Introdução}

O \textbf{$\lambda$-cálculo} é a linguagem de programação mais simples
que ainda captura a essência da computação por funções. Formulado por
Alonzo Church na década de 1930, ele consiste em apenas três
construções: variáveis, abstrações (funções anônimas) e aplicações
(chamadas de função). Apesar dessa aparente simplicidade, o cálculo
lambda não tipado é Turing-completo, qualquer função computável
pode ser expressa nele.

Neste capítulo, estudamos o $\lambda$-cálculo em duas formas. Primeiro,
apresentamos a versão \textbf{não tipada}, definindo sua sintaxe e as
semânticas big-step e small-step na estratégia \emph{call-by-value}.
Em seguida, o estendemos com valores booleanos e um sistema de tipos
simples (o \textbf{$\lambda$-cálculo simplesmente tipado}) e
demonstramos que esse sistema satisfaz as propriedades de progresso e
preservação. Por fim, apresentamos a implementação em Haskell do
interpretador para a versão não tipada.

\section{O $\lambda$-cálculo Não Tipado}

\subsection{Sintaxe}

A gramática do $\lambda$-cálculo não tipado tem três produções:

\[
  \begin{array}{rcll}
    t & ::= & x & \text{variável} \\
    & |   & \lambda x.\, t & \text{abstração} \\
    & |   & t\, t   & \text{aplicação}
\end{array} \]

Usamos as letras \(x, y, z, \ldots\) para variáveis e \(t, s, u,
\ldots\) para termos genéricos. Duas convenções de precedência tornam
a notação menos carregada de parênteses:

\begin{itemize}
  \item A \textbf{aplicação} é associativa à esquerda: \(t_1\, t_2\,
    t_3\) significa \((t_1\, t_2)\, t_3\).
  \item O \textbf{corpo de uma abstração} se estende o máximo possível à
    direita: \(\lambda x.\, t_1\, t_2\) significa \(\lambda
    x.\,(t_1\, t_2)\), não \((\lambda x.\, t_1)\, t_2\).
\end{itemize}

Abstrações com múltiplos parâmetros são abreviadas: \(\lambda x\,
y.\, t\) significa \(\lambda x.\,\lambda y.\, t\).

Uma \textbf{ocorrência livre} de \(x\) em \(t\) é uma ocorrência que
não está no escopo de nenhuma abstração \(\lambda x\). Um termo sem
variáveis livres é dito \textbf{fechado} (também chamado de
\emph{combinador}). O subconjunto de termos fechados é o que
consideraremos para avaliação.

\subsection{Valores}

Na estratégia call-by-value, os únicos \textbf{valores} são as
abstrações, que são as formas normais do cálculo. Variáveis e
aplicações só são valores se toda redução possível já foi realizada,
o que para termos fechados equivale a ser uma abstração:

\[
  \begin{array}{rcll} v & ::= & \lambda x.\, t & \text{(valor)}
  \end{array}
\]

\subsection{Semântica Big-Step Call-by-Value}

A semântica big-step define uma relação \(t \Downarrow v\) que
expressa que o termo \(t\) \textbf{avalia} ao valor \(v\) em um
número finito de passos. A estratégia call-by-value impõe que o
argumento de uma aplicação seja reduzido a um valor \textbf{antes} de
realizar a substituição (ao contrário de call-by-name, que substitui
o argumento sem avaliá-lo).

\[ \dfrac{}{\lambda x.\, t\ \Downarrow\ \lambda x.\, t} \quad\text{(B-Lam)}
\]

\[ \dfrac{ t_1 \Downarrow \lambda x.\, t \qquad t_2 \Downarrow v_2 \qquad t[x
\mapsto v_2] \Downarrow v }{ t_1\, t_2\ \Downarrow\ v } \quad\text{(B-App)} \]

A regra (B-Lam) expressa que uma abstração já é um valor: ela avalia
a si mesma. A regra (B-App) descreve os três passos de uma chamada de
função: avalia a função \(t_1\) até obter uma abstração \(\lambda
x.\, t\), avalia o argumento \(t_2\) até obter um valor \(v_2\), e
então avalia o corpo \(t\) com \(x\) substituído por \(v_2\).

A \textbf{substituição} \(t[x \mapsto v]\) é definida recursivamente
sobre a estrutura de \(t\):

\[
  \begin{array}{rcl}
    x[x \mapsto v] & = & v \\
    y[x \mapsto v] & = & y \quad (y \neq x) \\
    (\lambda x.\, t)[x \mapsto v] & = & \lambda x.\, t \\
    (\lambda y.\, t)[x \mapsto v] & = & \lambda y.\, (t[x \mapsto v]) \quad (y
    \neq x,\; y \notin \mathrm{fv}(v)) \\
    (t_1\, t_2)[x \mapsto v] & = & (t_1[x
    \mapsto v])\,(t_2[x \mapsto v])
\end{array} \]

A condição \(y \notin \mathrm{fv}(v)\) na regra para abstrações evita
a \textbf{captura de variável}: se \(y\) ocorre livre em \(v\), é
necessário primeiro renomear \(y\) para uma variável fresca
(\(\alpha\)-conversão).

Como exemplo, considere a avaliação do combinador identidade aplicado
a si mesmo:

\[ \dfrac{ \lambda x.\, x \Downarrow \lambda x.\, x \qquad \lambda x.\, x
    \Downarrow \lambda x.\, x \qquad (\lambda x.\, x)[x \mapsto
    \lambda x.\, x] =
  \lambda x.\, x \Downarrow \lambda x.\, x }{ (\lambda x.\, x)\,(\lambda x.\,
x) \Downarrow \lambda x.\, x } \quad\text{(B-App)} \]

\subsection{Semântica Small-Step Call-by-Value}

A semântica small-step define uma relação \(t \to t'\) que decompõe a
avaliação em passos elementares. É composta por uma \textbf{regra de
computação} e duas \textbf{regras de congruência}:

\textbf{Regra de computação (\(\beta\)-redução call-by-value):}

\[ \dfrac{}{(\lambda x.\, t_1)\, v_2 \to t_1[x \mapsto v_2]}
\quad\text{(E-Beta)} \]

A regra (E-Beta) só se aplica quando o argumento \(v_2\) já é um
valor, o que constitui a essência da estratégia call-by-value.

\textbf{Regras de congruência:}

\[ \dfrac{t_1 \to t_1'}{t_1\, t_2 \to t_1'\, t_2} \quad\text{(E-App1)} \qquad
\dfrac{t_2 \to t_2'}{v_1\, t_2 \to v_1\, t_2'} \quad\text{(E-App2)} \]

As regras de congruência determinam a \textbf{ordem de avaliação}:
(E-App1) reduz a função antes de olhar para o argumento; (E-App2) só
é aplicável quando a função já é um valor \(v_1\), e reduz o
argumento. Somente depois que ambos forem valores é que (E-Beta) pode disparar.

A sequência de reduções para o mesmo exemplo anterior é:

\[ (\lambda x.\, x)\,(\lambda x.\, x) \;\xrightarrow{\text{E-Beta}}\;
\lambda x.\, x \]

Como ambos os subtermos já são valores, a única regra aplicável é
(E-Beta) diretamente.

Um exemplo com mais passos, usando o combinador de aplicação
\(K = \lambda x.\,\lambda y.\,
x\):

\[ K\,(\lambda z.\,z)\,(\lambda w.\,w) \;\xrightarrow{\text{E-App1,
  E-Beta}}\; (\lambda y.\, \lambda z.\,z)\,(\lambda w.\,w)
\;\xrightarrow{\text{E-Beta}}\; \lambda z.\,z \]

Um termo é dito \textbf{preso} (\emph{stuck}) quando não é um
valor e nenhuma regra de redução se aplica a ele. No $\lambda$-cálculo
não tipado, uma variável livre é um exemplo de termo preso. A
versão tipada que estudamos a seguir elimina essa possibilidade para
termos bem tipados.

\section{$\lambda$-cálculo Simplesmente Tipado com Booleanos}

Estendemos o $\lambda$-cálculo com valores booleanos e um sistema de
tipos simples. Esta extensão, conhecida como \textbf{STLC+Bool}
(Simply Typed Lambda Calculus with Booleans), é um dos exemplos
fundamentais na teoria de tipos.

\subsection{Sintaxe Estendida}

Os tipos são:

\[
  \begin{array}{rcll}
    T & ::= & \mathbf{Bool} & \text{tipo booleano} \\
    & |   & T \to T & \text{tipo funcional}
\end{array} \]

Os termos são estendidos com os três construtores booleanos, e as
abstrações passam a carregar uma anotação de tipo no parâmetro:

\[
  \begin{array}{rcll} t & ::= & x & \text{variável} \\ & | & \lambda x {:}
    T.\, t & \text{abstração tipada} \\ & | & t\ t & \text{aplicação} \\ & | &
    \mathbf{true} & \text{constante verdadeiro} \\ & | & \mathbf{false} &
    \text{constante falso} \\ & | & \mathbf{if}\, t\, \mathbf{then}\, t\,
    \mathbf{else}\, t & \text{condicional}
\end{array} \]

Os valores são:

\[
  \begin{array}{rcll} v & ::= & \lambda x {:} T.\, t & \text{abstração} \\
    & | & \mathbf{true} & \\ & | & \mathbf{false} &
\end{array} \]

\subsection{Sistema de Tipos}

O sistema de tipos é definido por um \textbf{contexto} \(\Gamma\),
que é uma função parcial de variáveis para tipos, e pela relação de
tipagem \(\Gamma \vdash t : T\), lida como
``no contexto \(\Gamma\), o termo \(t\) tem tipo \(T\)''. As seis
regras de inferência são:

\[ \dfrac{}{\Gamma \vdash \mathbf{true} : \mathbf{Bool}} \;\text{(T-True)}
  \qquad \dfrac{}{\Gamma \vdash \mathbf{false} : \mathbf{Bool}}
\;\text{(T-False)} \]

\[ \dfrac{x : T \in \Gamma}{\Gamma \vdash x : T} \;\text{(T-Var)} \]

\[ \dfrac{\Gamma,\, x : T_1 \vdash t : T_2} {\Gamma \vdash \lambda x {:}
T_1.\, t : T_1 \to T_2} \;\text{(T-Abs)} \]

\[ \dfrac{\Gamma \vdash t_1 : T_1 \to T_2 \qquad \Gamma \vdash t_2 : T_1}
{\Gamma \vdash t_1\, t_2 : T_2} \;\text{(T-App)} \]

\[ \dfrac{ \Gamma \vdash t_1 : \mathbf{Bool} \qquad \Gamma \vdash t_2 : T
  \qquad \Gamma \vdash t_3 : T }{ \Gamma \vdash \mathbf{if}\ t_1\ \mathbf{then}\
t_2\ \mathbf{else}\ t_3 : T } \;\text{(T-If)} \]

A notação \(\Gamma, x : T\) denota o contexto obtido estendendo
\(\Gamma\) com a ligação \(x : T\). Escrevemos \(\vdash t : T\) (sem
\(\Gamma\)) quando o contexto é vazio; neste caso, \(t\) é um termo fechado.

Dois exemplos de derivações. O combinador identidade para booleanos
\(\lambda x \{:\} \mathbf{Bool}., x\) tem
tipo \(\mathbf{Bool} \to \mathbf{Bool}\):

\[ \dfrac{ \dfrac{x : \mathbf{Bool} \in \{x : \mathbf{Bool}\}} {x :
  \mathbf{Bool} \vdash x : \mathbf{Bool}} \;\text{(T-Var)} }{ \vdash \lambda x
  {:} \mathbf{Bool}.\, x : \mathbf{Bool} \to \mathbf{Bool} } \;\text{(T-Abs)}
\]

A função \(\lambda f \{:\} \mathbf{Bool} \to
  \mathbf{Bool}.\, f\,
\mathbf{true}\) tem tipo \((\mathbf{Bool}
\to \mathbf{Bool}) \to \mathbf{Bool}\):

\[ \dfrac{ \dfrac{f : \mathbf{Bool} \to \mathbf{Bool} \in \Gamma'} {\Gamma'
    \vdash f : \mathbf{Bool} \to \mathbf{Bool}} \quad \dfrac{}{\Gamma' \vdash
  \mathbf{true} : \mathbf{Bool}} }{ \vdash \lambda f {:} \mathbf{Bool} \to
    \mathbf{Bool}.\, f\, \mathbf{true} : (\mathbf{Bool} \to \mathbf{Bool}) \to
\mathbf{Bool} } \]

onde \(\Gamma' = f : \mathbf{Bool} \to \mathbf{Bool}\).

\subsection{Semântica Small-Step}

A semântica small-step do STLC+Bool estende as regras do cálculo não
tipado com as regras para condicionais:

\[ \dfrac{}{(\lambda x {:} T.\, t_1)\, v_2 \to t_1[x \mapsto v_2]}
\;\text{(E-Beta)} \]

\[ \dfrac{t_1 \to t_1'}{t_1\, t_2 \to t_1'\, t_2} \;\text{(E-App1)} \qquad
\dfrac{t_2 \to t_2'}{v_1\, t_2 \to v_1\, t_2'} \;\text{(E-App2)} \]

\[ \dfrac{}{\mathbf{if}\ \mathbf{true}\ \mathbf{then}\ t_2\ \mathbf{else}\ t_3
\to t_2} \;\text{(E-IfTrue)} \]

\[ \dfrac{}{\mathbf{if}\ \mathbf{false}\ \mathbf{then}\ t_2\ \mathbf{else}\ t_3
\to t_3} \;\text{(E-IfFalse)} \]

\[ \dfrac{t_1 \to t_1'} {\mathbf{if}\ t_1\ \mathbf{then}\ t_2\ \mathbf{else}\
  t_3 \to \mathbf{if}\ t_1'\ \mathbf{then}\ t_2\ \mathbf{else}\ t_3}
\;\text{(E-If)} \]

\subsection{Formas Canônicas}

Antes das demonstrações de progresso e preservação, precisamos de um
lema que relaciona valores a tipos. O \textbf{lema das formas
canônicas} afirma que os únicos valores de cada tipo têm a forma esperada.

\textbf{Lema (Formas Canônicas).} Seja \(v\) um valor com \(\vdash v : T\).

\begin{enumerate}
  \item Se \(T = \mathbf{Bool}\), então \(v = \mathbf{true}\) ou
    \(v = \mathbf{false}\).
  \item Se \(T = T_1 \to T_2\), então \(v = \lambda x {:} T_1.\, t\) para
    algum \(x\) e \(t\).
\end{enumerate}

\emph{Demonstração.} Direta por inspeção das regras de tipagem. Os
únicos valores são \(\mathbf{true}\), \(\mathbf{false}\) e \(\lambda
x {:} T_1.\, t\). A regra (T-True) atribui tipo \(\mathbf{Bool}\) a
\(\mathbf{true}\); a regra (T-False) atribui tipo \(\mathbf{Bool}\) a
\(\mathbf{false}\); a regra (T-Abs) atribui tipo \(T_1 \to T_2\) a
\(\lambda x {:} T_1.\, t\). Nenhum outro valor pode ter
\(\mathbf{Bool}\) ou um tipo funcional.

\subsection{Progresso}

O teorema de progresso garante que um programa bem tipado fechado
nunca fica preso: ele ou já é um valor, ou pode dar mais um passo
de redução.

\textbf{Teorema (Progresso).} Se \(\vdash t : T\) (com contexto
vazio), então \(t\) é um valor ou \(\exists\, t'\) tal que \(t \to t'\).

\emph{Demonstração.} Por indução estrutural sobre a derivação de
\(\vdash t : T\).

\textbf{Caso (T-True):} \(t = \mathbf{true}\) é um valor.

\textbf{Caso (T-False):} \(t = \mathbf{false}\) é um valor.

\textbf{Caso (T-Var):} \(t = x\). Este caso não pode ocorrer, pois
\(\Gamma = \emptyset\) implica que nenhuma
variável tem tipo definido.

\textbf{Caso (T-Abs):} \(t = \lambda x {:} T_1.\, t_1\) é um valor.

\textbf{Caso (T-App):} \(t = t_1\, t_2\), com \(\vdash t_1 : T_1 \to
T_2\) e \(\vdash t_2 : T_1\). Pela hipótese de indução aplicada a \(t_1\):

\begin{itemize}
  \item
    \emph{Subcaso \(t_1 \to t_1'\):} Por (E-App1), temos
    \(t\_1\, t\_2 \to t\_1', t\_2\).
  \item
    \emph{Subcaso \(t_1\) é um valor:} Pela hipótese de indução
    aplicada a \(t_2\):

    \begin{itemize}
      \item
        \emph{Subcaso \(t_2 \to t_2'\):} Por (E-App2), temos
        \(t\_1\, t\_2 \to
        t\_1\, t\_2'\).
      \item
        \emph{Subcaso \(t_2\) é um valor:} Como \(t_1\) é um valor
        com tipo \(T\_1 \to T\_2\), o Lema
        das Formas Canônicas (item 2) garante que \(t\_1
        = \lambda x \{:\} T\_1.\, t\).
        Logo, por (E-Beta): \[ (\lambda x {:} T_1.\,
        t)\, t_2 \to t[x \mapsto t_2] \]
    \end{itemize}
\end{itemize}

Em todos os subcasos, \(t\) dá um passo de redução.

\textbf{Caso (T-If):} \(t =
  \mathbf{if}~t\_1~\mathbf{then}~t\_2~\mathbf{else}\\
t\_3\), com \(\vdash t_1 : \mathbf{Bool}\). Pela
hipótese de indução aplicada a \(t_1\):

\begin{itemize}
  \item
    \emph{Subcaso \(t\_1 \to t\_1'\):} Por
    (E-If): \[ \mathbf{if}\ t_1\ \mathbf{then}\
      t_2\ \mathbf{else}\ t_3 \;\to\; \mathbf{if}\ t_1'\ \mathbf{then}\ t_2\
    \mathbf{else}\ t_3 \]
  \item
    \emph{Subcaso \(t_1\) é um valor:} Como \(t_1 : \mathbf{Bool}\) é
    um valor, o Lema das Formas Canônicas (item 1) garante que \(t_1
    = \mathbf{true}\) ou \(t_1 = \mathbf{false}\). Em cada caso,
    (E-IfTrue) ou (E-IfFalse) se aplica.
\end{itemize}

Em todos os subcasos, \(t\) dá um passo de redução.

\subsection{Preservação}

O teorema de preservação garante que a redução não altera o tipo do
termo: se um termo bem tipado dá um passo, o resultado continua bem
tipado com o mesmo tipo.

Antes da demonstração, precisamos de um lema técnico fundamental
sobre substituição.

\textbf{Lema (Substituição).} Se \(\Gamma, x : T_1 \vdash t : T_2\) e
\(\Gamma \vdash v : T\_1\), então \(\Gamma
\vdash t[x \mapsto v] : T_2\).

\emph{Demonstração (esboço).} Por indução estrutural sobre \(t\).

\begin{itemize}
  \item
    \emph{\(t = x\):} A substituição produz \(v\), e
    \(\Gamma \vdash v : T\_1 = T\_2\).
  \item
    \emph{\(t = y\) com \(y \neq x\):} A substituição não altera
    \(y\), e \(y : T\_2 \in \Gamma\) ainda vale.
  \item
    \emph{\(t = \lambda y {:} T.\, t'\) com \(y \neq x\):} Por
    hipótese, \(\Gamma, x : T\_1, y : T \vdash t' :
    T\_2'\). Por hipótese de indução,
    \(\Gamma, y : T \vdash t'{[}x \mapsto v{]} :
    T\_2'\). Logo, pela regra (T-Abs),
    \(\Gamma \vdash \lambda y \{:\} T.\,
    t'{[}x \mapsto v{]} : T \to T\_2'\).
  \item
    \emph{\(t = t_a\, t_b\):} Por hipótese, \(\Gamma, x
    : T\_1 \vdash t\_a : T \to T\_2\) e \(\Gamma, x :
    T_1 \vdash t_b : T\). Por hipótese de indução em ambos: \(\Gamma
    \vdash t_a[x \mapsto v] : T \to T_2\) e \(\Gamma
    \vdash t\_b{[}x \mapsto v{]} : T\). Por (T-App):
    \(\Gamma \vdash (t\_a, t\_b){[}x \mapsto v{]} :
    T\_2\).
  \item
    \emph{\(t = \mathbf{true}\), \(t = \mathbf{false}\),
    \(t = \mathbf{if}\ldots\):} Análogos,
    usando (T-True), (T-False) e (T-If).
\end{itemize}

\textbf{Teorema (Preservação).} Se \(\Gamma \vdash t : T\) e \(t \to
t'\), então \(\Gamma \vdash t' : T\).

\emph{Demonstração.} Por indução sobre a derivação de \(t \to t'\).

\textbf{Caso (E-Beta):} \(t = (\lambda x \{:\}
  T\_1.\, t\_1)\, v\_2 \to t\_1{[}x \mapsto
v\_2{]} = t'\).

A derivação de \(\Gamma \vdash t : T_2\) termina com (T-App), portanto:

\begin{itemize}
  \item
    \(\Gamma \vdash \lambda x {:} T_1.\, t_1 : T_1 \to T_2\),
    derivada por (T-Abs), donde \(\Gamma, x : T_1 \vdash t_1 : T_2\).
  \item
    \(\Gamma \vdash v_2 : T_1\).
\end{itemize}

Pelo Lema da Substituição: \(\Gamma \vdash t_1[x \mapsto v_2] : T_2\).

\textbf{Caso (E-App1):} \(t = t_1\, t_2 \to t_1'\, t_2 = t'\), com
\(t\_1 \to t\_1'\).

A derivação de \(\Gamma \vdash t : T_2\) termina com (T-App):

\begin{itemize}
  \item
    \(\Gamma \vdash t_1 : T_1 \to T_2\)
  \item
    \(\Gamma \vdash t_2 : T_1\)
\end{itemize}

Por hipótese de indução: \(\Gamma \vdash t_1' : T_1 \to T_2\). Por
(T-App): \(\Gamma \vdash t_1'\, t_2 : T_2\).

\textbf{Caso (E-App2):} \(t = v_1\, t_2 \to v_1\, t_2' = t'\), com
\(t\_2 \to t\_2'\).

Analogamente, (T-App) fornece \(\Gamma \vdash v_1 : T_1 \to T_2\) e
\(\Gamma \vdash t\_2 : T\_1\). Por hipótese
de indução: \(\Gamma \vdash t_2' : T_1\). Por (T-App): \(\Gamma
\vdash v_1\, t_2' : T_2\).

\textbf{Caso (E-IfTrue):} \(t =
  \mathbf{if}~\mathbf{true}~\mathbf{then}~t\_2\\
\mathbf{else}~t\_3 \to t\_2 = t'\).

A derivação termina com (T-If), donde \(\Gamma \vdash t_2 : T\).
Portanto \(\Gamma \vdash t' : T\).

\textbf{Caso (E-IfFalse):} \(t =
  \mathbf{if}~\mathbf{false}~\mathbf{then}~t\_2\\
\mathbf{else}~t\_3 \to t\_3 = t'\).

Análogo: (T-If) fornece \(\Gamma \vdash t_3 : T\).

\textbf{Caso (E-If):} \(t =
  \mathbf{if}~t\_1~\mathbf{then}~t\_2~\mathbf{else}~t\_3 \to
  \mathbf{if}~t\_1'~\mathbf{then}~t\_2~\mathbf{else}~t\_3 =
t'\), com \(t_1 \to t_1'\).

A derivação termina com (T-If):

\begin{itemize}
  \item
    \(\Gamma \vdash t_1 : \mathbf{Bool}\)
  \item
    \(\Gamma \vdash t_2 : T\)
  \item
    \(\Gamma \vdash t_3 : T\)
\end{itemize}

Por hipótese de indução: \(\Gamma \vdash t_1' : \mathbf{Bool}\). Por
(T-If): \(\Gamma \vdash
\mathbf{if}~t\_1'~\mathbf{then}~t\_2~\mathbf{else}~t\_3 : T\).

Os dois teoremas juntos estabelecem a \textbf{correção do tipo} do
STLC+Bool: toda sequência de reduções a partir de um termo fechado e
bem tipado produz, a cada passo, termos bem tipados do mesmo tipo, e
jamais produz um termo preso.

\section{Closure Conversion}

\subsection{Motivação}

O $\lambda$-cálculo trata funções como valores de primeira classe: uma
abstração pode ser passada como argumento, retornada como resultado e
armazenada em estruturas de dados. Em uma linguagem de baixo nível
como C ou em uma representação intermediária como a IRT, toda função
é um procedimento com endereço estático; não há, nativamente, o
conceito de um valor funcional que carrega variáveis livres.

A \textbf{closure conversion} (conversão de closures) é a
transformação que torna esse gap explícito e o resolve. Ela separa
cada abstração em dois componentes:

\begin{itemize}
  \item Um \textbf{código promovido} (lifted): uma função de nível superior,
    completamente fechada (sem variáveis livres), que recebe o ambiente
    capturado como parâmetro explícito.
  \item Um \textbf{registro de closure}: uma estrutura de dados criada no
    ponto em que a abstração aparece no programa-fonte, contendo uma
    referência ao código promovido e os valores das variáveis
    capturadas.
\end{itemize}

\noindent
O resultado da closure conversion é um programa onde toda função é
fechada e toda aplicação passa o ambiente capturado explicitamente.
Nesta seção formalizamos essa transformação para o $\lambda$-cálculo não
tipado, usando a linguagem TImp como alvo.

\subsection{Variáveis Livres}

A análise de variáveis livres é a pré-condição da closure conversion.
Definimos \(\mathrm{FV}(t)\), o conjunto de variáveis livres de um
termo \(t\), por indução estrutural:

\[
  \begin{array}{rcl}
    \mathrm{FV}(x)        & = & \{x\} \\
    \mathrm{FV}(\lambda x.\, t) & = & \mathrm{FV}(t) \setminus \{x\} \\
    \mathrm{FV}(t_1\, t_2) & = & \mathrm{FV}(t_1) \cup \mathrm{FV}(t_2)
  \end{array}
\]

Um termo \(t\) é \textbf{fechado} quando \(\mathrm{FV}(t) = \emptyset\).
Assumimos que o termo de entrada da closure conversion é fechado;
variáveis livres restantes correspondem a erros do programador.

\subsection{Anotação de Lambdas}

O primeiro passo é percorrer o termo e atribuir a cada abstração um
\textbf{identificador único} \(k \in \mathbb{N}\) e registrar a lista
das variáveis que ela captura. Definimos os termos anotados
\(\hat{t}\) por:

\[
  \hat{t} \;::=\; \hat{x}
  \;\mid\; \hat{t}\; \hat{t}
  \;\mid\; \lambda^k\, x\,[\overline{y}]\,.\;\hat{t}
\]

onde \(\overline{y} = [y_1, \ldots, y_j]\) é a lista (em alguma ordem
fixa) das variáveis livres da subclosure, isto é,
\(\overline{y} = \mathrm{FV}(\hat{t}) \setminus \{x\}\).

O conjunto de variáveis livres de um termo anotado é definido por:

\[
  \begin{array}{rcl}
    \mathrm{FV}(\hat{x}) & = & \{x\} \\
    \mathrm{FV}(\hat{t}_1\; \hat{t}_2) & = & \mathrm{FV}(\hat{t}_1) \cup \mathrm{FV}(\hat{t}_2) \\
    \mathrm{FV}(\lambda^k\, x\,[\overline{y}]\,.\;\hat{t}) & = & \{\overline{y}\}
  \end{array}
\]

A função de anotação \(\mathit{ann}\) percorre o termo em
pós-ordem (abstrações internas recebem IDs menores):

\[
  \begin{array}{rcl}
    \mathit{ann}(x) & = & \hat{x} \\[4pt]
    \mathit{ann}(t_1\, t_2) & = & \mathit{ann}(t_1)\;\mathit{ann}(t_2) \\[4pt]
    \mathit{ann}(\lambda x.\, t) & = &
      \lambda^k\, x\,[\overline{y}]\,.\;\hat{t}
      \quad \text{onde } \hat{t} = \mathit{ann}(t),\;
      k \text{ fresco},\;
      \overline{y} = \mathrm{FV}(\hat{t}) \setminus \{x\}
  \end{array}
\]

\subsection{O Tipo Closure}

A closure conversion utiliza um \textbf{registro uniforme} para
representar toda função como valor em tempo de execução. Seja
\(N = \max\{|\overline{y}| : \lambda^k\, x\,[\overline{y}]\,.\;\hat{t}
\text{ ocorre em } \mathit{ann}(t)\}\) o número máximo de variáveis
capturadas por qualquer abstração do programa. O tipo alvo é:

\[
  \mathbf{record}\;\mathtt{Closure}\;\{
    \;\mathtt{tag} : \mathbf{int};\;
    \mathtt{f}_0 : \mathtt{Closure};\;
    \ldots;\;
    \mathtt{f}_{N-1} : \mathtt{Closure}\;\}
\]

O campo \(\mathtt{tag}\) identifica qual lambda a closure representa
(o identificador \(k\)), e os campos \(\mathtt{f}_i\) guardam os
valores das variáveis capturadas. Os campos não utilizados por uma
closure particular simplesmente não são inicializados.

\subsection{A Tradução}

A tradução \(\llbracket \hat{t}\rrbracket_\sigma\) converte um termo
anotado em uma expressão TImp, dado um \textbf{ambiente de
substituição} \(\sigma : \mathit{Var} \rightharpoonup
\mathit{ExpTImp}\) que mapeia variáveis-fonte para expressões-alvo.

\[
  \llbracket \hat{x} \rrbracket_\sigma \;=\;
  \begin{cases}
    \sigma(x) & \text{se } x \in \mathrm{dom}(\sigma) \\
    x         & \text{caso contrário}
  \end{cases}
\]

\[
  \llbracket \hat{t}_1\; \hat{t}_2 \rrbracket_\sigma \;=\;
  \mathtt{apply}\bigl(\llbracket \hat{t}_1 \rrbracket_\sigma,\;
  \llbracket \hat{t}_2 \rrbracket_\sigma\bigr)
\]

\[
  \llbracket \lambda^k\, x\,[y_1,\ldots,y_j]\,.\;\hat{t}
  \rrbracket_\sigma \;=\;
  \mathbf{new}\;\mathtt{Closure}\;\bigl\{
    \mathtt{tag} = k,\;
    \mathtt{f}_0 = \sigma(y_1),\;
    \ldots,\;
    \mathtt{f}_{j-1} = \sigma(y_j)
  \bigr\}
\]

A regra para uma abstração \emph{não} compila o corpo neste ponto: ela
apenas cria o registro de closure com a identificação do código
promovido (\(\mathtt{tag}\)) e os valores capturados. O corpo é
compilado separadamente, na geração das funções promovidas.

\subsection{Funções Promovidas}

Para cada lambda anotado \((k, x, [y_1,\ldots,y_j], \hat{t})\) coletado
do termo, gera-se a função TImp:

\[
  \mathbf{fn}\;\mathtt{lam}\_k\;
  (\mathtt{\_env} : \mathtt{Closure},\;
  \mathtt{\_arg} : \mathtt{Closure}) :
  \mathtt{Closure}\;\{
  \;\mathbf{return}\;
  \llbracket \hat{t} \rrbracket_{\sigma_k}\;\}
\]

onde o ambiente de substituição \(\sigma_k\) é:

\[
  \sigma_k \;=\;
  \bigl[y_1 \mapsto \mathtt{\_env.f}_0,\;\ldots,\;
  y_j \mapsto \mathtt{\_env.f}_{j-1},\;
  x \mapsto \mathtt{\_arg}\bigr]
\]

A substituição em \(\sigma_k\) mapeia cada variável capturada ao
campo correspondente de \texttt{\_env} e o parâmetro formal \(x\) ao
argumento recebido. Toda chamada recursiva a \(\llbracket \cdot
\rrbracket\) dentro do corpo usa \(\sigma_k\), de modo que o corpo
inteiro é compilado como uma única expressão TImp.

\subsection{Despacho Dinâmico}

Como o campo \texttt{tag} de uma closure determina qual função promovida
chamar, é necessária uma função de \textbf{despacho} que, dado um par
(closure, argumento), redirecione para a função promovida correta. Com
\(n\) lambdas de tags \(k_1 < k_2 < \cdots < k_n\):

\[
  \mathbf{fn}\;\mathtt{apply}\;
  (\mathtt{\_f} : \mathtt{Closure},\;
  \mathtt{\_arg} : \mathtt{Closure}) :
  \mathtt{Closure}\;\{
\]
\[
  \quad \mathbf{if}\;\mathtt{\_f.tag} = k_1\;
  \mathbf{then}\;\{\mathbf{return}\;\mathtt{lam}_{k_1}(\mathtt{\_f}, \mathtt{\_arg})\;\}
\]
\[
  \quad \mathbf{else}\;\{\;\mathbf{if}\;\mathtt{\_f.tag} = k_2\;\mathbf{then}\;
  \{\mathbf{return}\;\mathtt{lam}_{k_2}(\mathtt{\_f}, \mathtt{\_arg})\;\}
\]
\[
  \quad \ldots
\]
\[
  \quad\mathbf{else}\;\{\mathbf{return}\;
  \mathbf{new}\;\mathtt{Closure}\;\{\mathtt{tag} = {-1}\}\;\}\;\}\;\}
\]

O bloco \(\mathbf{else}\) final é inalcançável para programas bem
formados, pois todo valor-closure criado pela tradução tem um tag
válido.

\subsection{Correção da Tradução}

A closure conversion preserva a semântica do lambda cálculo: para
qualquer termo fechado \(t\) tal que \(t \Downarrow v\) na semântica
big-step call-by-value, a execução do programa TImp gerado termina
imprimindo a closure correspondente ao valor \(v\). Mais precisamente,
o valor \(v = \lambda x.\, t'\) é representado em tempo de execução
pela closure \(\mathbf{new}\;\mathtt{Closure}\;\{\mathtt{tag} =
k,\;\ldots\}\) onde \(k\) é o identificador atribuído à abstração
\(\lambda x.\, t'\) e os campos \(\mathtt{f}_i\) carregam as
representações dos valores de suas variáveis livres.

A invariante central é: \textbf{toda closure em tempo de execução
representa exatamente um valor do $\lambda$-cálculo}. Os campos
\(\mathtt{f}_i\) são sempre closures válidas (nunca campos não
inicializados são lidos), porque a tradução \(\llbracket \lambda^k\,
x\,[y_1,\ldots,y_j]\,.\;\hat{t}\rrbracket_\sigma\) escreve
exatamente os \(j\) campos usados por \(\mathtt{lam}_k\).

\subsection{Exemplo Trabalhado}

Consideremos o combinador K aplicado ao combinador identidade:

\[
  (\lambda f.\,\lambda x.\, f)\;(\lambda y.\, y)
\]

\paragraph{Passo 1: Anotação.}

Percorrendo em pós-ordem:
\begin{itemize}
  \item \(\lambda y.\, y\) tem \(\mathrm{FV} = \emptyset\); recebe tag \(0\).
  \item \(\lambda x.\, f\) tem \(\mathrm{FV}(\hat{t}) \setminus \{x\} = \{f\}\); recebe tag \(1\).
  \item \(\lambda f.\,\ldots\) tem \(\mathrm{FV} = \emptyset\); recebe tag \(2\).
\end{itemize}

O termo anotado é:
\[
  \bigl(\lambda^2 f\,[]\,.\;\lambda^1 x\,[f]\,.\;\hat{f}\bigr)\;
  \bigl(\lambda^0 y\,[]\,.\;\hat{y}\bigr)
\]

\paragraph{Passo 2: Tipo Closure.}

O máximo de variáveis capturadas é 1 (em \(\lambda^1\)). Logo:
\[
  \mathbf{record}\;\mathtt{Closure}\;\{\;\mathtt{tag}: \mathbf{int};\;\mathtt{f}_0: \mathtt{Closure}\;\}
\]

\paragraph{Passo 3: Funções promovidas.}

Para \(\lambda^0 y\,[]\,.\;\hat{y}\) com \(\sigma_0 = [y \mapsto \mathtt{\_arg}]\):
\begin{verbatim}
fn lam_0(_env: Closure, _arg: Closure) : Closure {
  return _arg;
}
\end{verbatim}

Para \(\lambda^1 x\,[f]\,.\;\hat{f}\) com \(\sigma_1 = [f \mapsto \mathtt{\_env.f}_0,\, x \mapsto \mathtt{\_arg}]\):
\begin{verbatim}
fn lam_1(_env: Closure, _arg: Closure) : Closure {
  return _env.f0;   /* [f -> _env.f0]_sigma1(hat{f}) */
}
\end{verbatim}

Para \(\lambda^2 f\,[]\,.\;\lambda^1 x\,[f]\,.\;\hat{f}\) com \(\sigma_2 = [f \mapsto \mathtt{\_arg}]\):
\[
  \llbracket \lambda^1 x\,[f]\,.\;\hat{f} \rrbracket_{\sigma_2}
  = \mathbf{new}\;\mathtt{Closure}\;\{\mathtt{tag}=1,\;
    \mathtt{f}_0 = \sigma_2(f)\}
  = \mathbf{new}\;\mathtt{Closure}\;\{\mathtt{tag}=1,\;
    \mathtt{f}_0 = \mathtt{\_arg}\}
\]
\begin{verbatim}
fn lam_2(_env: Closure, _arg: Closure) : Closure {
  return new Closure { tag = 1, f0 = _arg };
}
\end{verbatim}

\paragraph{Passo 4: Bloco principal.}

\[
  \llbracket (\lambda^2 f\,[]\,.\;\ldots)\;(\lambda^0 y\,[]\,.\;\hat{y}) \rrbracket_\emptyset
  = \mathtt{apply}\bigl(
      \mathbf{new}\;\mathtt{Closure}\{\mathtt{tag}=2\},\;
      \mathbf{new}\;\mathtt{Closure}\{\mathtt{tag}=0\}
    \bigr)
\]

\begin{verbatim}
var _res : Closure = apply(new Closure {tag=2}, new Closure {tag=0});
print _res;
\end{verbatim}

\paragraph{Execução.}

\begin{align*}
  &\mathtt{apply}(\{\mathtt{tag}=2\},\;\{\mathtt{tag}=0\})
  \\
  \to\;&\mathtt{lam\_2}(\{\mathtt{tag}=2\},\;\{\mathtt{tag}=0\})
  \\
  =\;&\mathbf{new}\;\mathtt{Closure}\;\{\mathtt{tag}=1,\;
       \mathtt{f}_0=\{\mathtt{tag}=0\}\}
\end{align*}

O resultado é a closure para \(\lambda x.\, f\) com \(f\) capturado
como \(\lambda y.\, y\). O bloco principal imprime:
\begin{verbatim}
Closure { tag = 1; f0 = Closure { tag = 0; }; }
\end{verbatim}

\section{Implementação em Haskell}

Apresentamos a implementação do interpretador para o \textbf{cálculo
lambda não tipado} em Haskell. A implementação segue a semântica
big-step call-by-value com um modelo de avaliação baseado em
\textbf{ambientes}, que substitui a operação de substituição textual
por uma associação de nomes a valores em tempo de execução.

\subsection{Representação de Termos}

A árvore de sintaxe abstrata reflete diretamente a gramática:

\begin{minted}{haskell}
type Name = String

data Term
  = Var Name
  | Lam Name Term
  | App Term Term
\end{minted}

\subsection{Valores e Ambientes}

Como valores são apenas abstrações, representamo-los como
\textbf{closures}, pares compostos pela abstração e pelo ambiente
léxico em que ela foi criada. O ambiente captura os valores das
variáveis livres da abstração no momento em que ela foi criada,
realizando o escopo léxico.

\begin{minted}{haskell}
data Value = VClosure Name Term Env

type Env = Map Name Value
\end{minted}

\subsection{O Interpretador}

A função \haskell{eval} implementa as regras (B-Lam)
e (B-App) diretamente. Para detectar potenciais não-terminações, um
contador de passos é mantido na mônada \haskell{State Int}:

\begin{minted}{haskell}
eval :: Env -> Term -> EvalM Value
eval env (Lam x body) =
  return (VClosure x body env)
eval env (Var x) =
  case Map.lookup x env of
    Just v  -> return v
    Nothing -> throwError $ "Unbound variable: " ++ x
eval env (App t1 t2) = do
    tick
    v1 <- eval env t1
    v2 <- eval env t2
    case v1 of
      VClosure x body closEnv ->
        eval (Map.insert x v2 closEnv) body
\end{minted}

A função \haskell{tick} incrementa o contador de
passos e interrompe a avaliação se o limite configurado (100 000
passos) for atingido, evitando que o REPL bloqueie indefinidamente em
termos como \((\lambda x.\, x\, x)\,(\lambda x.\, x\, x)\).

\subsection{A Interface REPL}

O executável \verb|lambda| oferece um REPL
interativo que lê um termo, avalia e imprime o resultado. A interface
aceita tanto \verb|\\| quanto
\verb|λ| como símbolo de abstração, e múltiplos
parâmetros em uma mesma abstração são abreviados:

\begin{verbatim}
Untyped Lambda Calculus
  syntax:  \x. t   |   t t   |   x
  multi-binder:  \x y z. t
  quit:    :q

λ> \x. x
--> λx. x

λ> (\x. x) (\y. y)
--> λy. y

λ> \x y. x
--> λx y. x

λ> (\x y. x) (\a. a) (\b. b)
--> λa. a

λ> (\x. x x) (\x. x x)
Error: Step limit exceeded after 100000 steps (possible non-terminating term)
\end{verbatim}

O modelo baseado em ambientes dispensa a operação de substituição
textual (com a correspondente necessidade de \(\alpha\)-conversão),
ao custo de que o resultado exibido mostra o corpo sintático da
closure, sem expandir as variáveis capturadas no ambiente. Para os
propósitos pedagógicos do REPL, esta representação é suficiente.

\subsection{Compilação para TImp via Closure Conversion}

A formalização da seção anterior se traduz diretamente em Haskell no
módulo \verb|Lambda.Backend.TImp.LambdaToTImp|. Apresentamos os
pontos de correspondência mais relevantes entre a teoria e o código.

\paragraph{Termos anotados.}

O tipo \haskell{Ann} representa os termos anotados \(\hat{t}\):

\begin{minted}{haskell}
type LamId = Int

data Ann
  = AVar Name
  | ALam LamId Name [Name] Ann  -- k, x, [y1,...,yj], corpo
  | AApp Ann Ann
\end{minted}

O construtor \haskell{ALam k x fvs body} corresponde a
\(\lambda^k\, x\,[\overline{y}]\,.\;\hat{t}\): \haskell{k} é o
identificador único, \haskell{fvs} é a lista \(\overline{y}\) de
variáveis capturadas e \haskell{body} é o corpo já anotado.

\paragraph{Variáveis livres de termos anotados.}

A função \haskell{freeVarsAnn} implementa a definição formal de
\(\mathrm{FV}\) sobre termos anotados:

\begin{minted}{haskell}
freeVarsAnn :: Ann -> Set Name
freeVarsAnn (AVar x)         = Set.singleton x
freeVarsAnn (AApp f a)       = freeVarsAnn f `Set.union` freeVarsAnn a
freeVarsAnn (ALam _ _ fvs _) = Set.fromList fvs
\end{minted}

Para um \haskell{ALam}, as variáveis livres são exatamente as já
armazenadas em \haskell{fvs}, sem necessidade de percorrer o corpo
novamente.

\paragraph{Passagem de anotação.}

A função \haskell{annotate} implementa \(\mathit{ann}\). O estado
monádico fornece identificadores frescos; o percurso é em pós-ordem,
de modo que lambdas internas recebem IDs menores:

\begin{minted}{haskell}
type AnnM = State Int

annotate :: Term -> AnnM Ann
annotate (Var x)   = return (AVar x)
annotate (App f a) = AApp <$> annotate f <*> annotate a
annotate (Lam x t) = do
  body <- annotate t
  k    <- freshId
  let fvs = Set.toList (Set.delete x (freeVarsAnn body))
  return (ALam k x fvs body)
\end{minted}

\paragraph{A tradução \(\llbracket \cdot \rrbracket_\sigma\).}

A função \haskell{compileExprS} implementa \(\llbracket \cdot
\rrbracket_\sigma\), recebendo o ambiente de substituição \(\sigma\)
como uma lista de pares \haskell{(Name, Exp)}:

\begin{minted}{haskell}
type Subst = [(Name, Exp)]

compileExprS :: Subst -> Ann -> Exp
compileExprS s (AVar x) =
  case lookup x s of { Just e -> e; Nothing -> EVar x }
compileExprS s (AApp f a) =
  ECall "apply" [compileExprS s f, compileExprS s a]
compileExprS s (ALam k _ fvs _) =
  ENew "Closure" $
    ("tag", EInt k) :
    [("f" ++ show i, compileExprS s (AVar v)) | (i,v) <- zip [0..] fvs]
\end{minted}

A equação para \haskell{ALam} implementa precisamente a regra:
\[
  \llbracket \lambda^k\, x\,[y_1,\ldots,y_j]\,.\;\hat{t}
  \rrbracket_\sigma
  = \mathbf{new}\;\mathtt{Closure}\;\{\mathtt{tag}=k,\;
  \mathtt{f}_0=\sigma(y_1),\;\ldots,\;\mathtt{f}_{j-1}=\sigma(y_j)\}
\]
O corpo \haskell{body} \emph{não} é compilado aqui: ele é compilado
separadamente, em \haskell{genLiftedFunc}.

\paragraph{Funções promovidas.}

A função \haskell{genLiftedFunc} gera cada \(\mathtt{lam}_k\). O
ambiente de substituição \(\sigma_k\) é construído explicitamente
como uma lista de pares antes de chamar \haskell{compileExprS}:

\begin{minted}{haskell}
genLiftedFunc :: (LamId, Name, [Name], Ann) -> FuncDecl
genLiftedFunc (k, x, fvs, body) =
  FuncDecl ("lam_" ++ show k)
    [Param "_env" (TRecord "Closure"), Param "_arg" (TRecord "Closure")]
    (RTTy (TRecord "Closure"))
    [SReturn (Just (compileExprS subst body))]
  where
    subst = [( v, EField (EVar "_env") ("f" ++ show i))
             | (i, v) <- zip [0..] fvs]
            ++ [(x, EVar "_arg")]
\end{minted}

O ambiente \haskell{subst} é exatamente \(\sigma_k = [y_i \mapsto
\mathtt{\_env.f}_{i-1}] \cup [x \mapsto \mathtt{\_arg}]\). O corpo
inteiro compila para uma única expressão TImp via
\haskell{compileExprS subst body}.

\paragraph{Despacho dinâmico.}

A função \haskell{genApplyFunc} constrói a cadeia de condicionais
aninhados. A lista \haskell{lids} contém os IDs de todos os lambdas:

\begin{minted}{haskell}
genApplyFunc :: [LamId] -> FuncDecl
genApplyFunc lids =
  FuncDecl "apply"
    [Param "_f" (TRecord "Closure"), Param "_arg" (TRecord "Closure")]
    (RTTy (TRecord "Closure"))
    [buildDispatch lids]
  where
    buildDispatch [] =
      SReturn (Just (ENew "Closure" [("tag", EInt (-1))]))
    buildDispatch (k:ks) =
      SIf (EField (EVar "_f") "tag" :=: EInt k)
          [SReturn (Just (ECall ("lam_" ++ show k) [EVar "_f", EVar "_arg"]))]
          [buildDispatch ks]
\end{minted}

\paragraph{Ponto de entrada.}

A função \haskell{lambdaToTImp} reúne todas as fases e produz um
programa TImp completo:

\begin{minted}{haskell}
lambdaToTImp :: Term -> TImp
lambdaToTImp term = TImp decls mainBlock
  where
    ann     = evalState (annotate term) 0
    lambdas = collectLambdas ann
    maxFVs  = maximum (0 : [length fvs | (_,_,fvs,_) <- lambdas])
    lids    = [k | (k,_,_,_) <- lambdas]
    decls   = DRecord (genClosureRecord maxFVs)
              : map (DFunc . genLiftedFunc) lambdas
              ++ [DFunc (genApplyFunc lids)]
    mainBlock =
      [ SDecl "_res" (TRecord "Closure") (compileExprS [] ann)
      , SPrint (EVar "_res") ]
\end{minted}

\paragraph{Verificação de tipos e inicialização parcial.}

O tipo \haskell{Closure} tem \(N\) campos \(\mathtt{f}_i\), mas uma
closure que captura \(j < N\) variáveis inicializa apenas
\(\mathtt{f}_0, \ldots, \mathtt{f}_{j-1}\). O verificador de tipos
de TImp foi ligeiramente relaxado para aceitar inicialização parcial
de registros: a regra (New) verifica que todos os campos
\emph{fornecidos} têm o tipo correto e que nenhum campo
\emph{extra} foi fornecido, mas não exige que todos os campos
declarados sejam inicializados. Em tempo de execução, os campos não
inicializados nunca são acessados para closures válidas, pois
\(\mathtt{lam}_k\) acessa apenas \(\mathtt{\_env.f}_0, \ldots,
\mathtt{\_env.f}_{j-1}\) e \(j\) é o número de campos que a closure
com tag \(k\) efetivamente inicializou.

\section{Conclusão}

O $\lambda$-cálculo é o núcleo de toda linguagem funcional moderna. As
suas semânticas big-step e small-step call-by-value definem de forma
precisa o que significa avaliar uma função: primeiro o operador,
depois o argumento, e só então a substituição.

A extensão com booleanos e tipos simples produziu o STLC+Bool, para o
qual demonstramos os dois teoremas centrais da teoria de tipos:
progresso e preservação. O progresso garante que programas bem
tipados nunca ficam presos; a preservação garante que a tipagem é
invariante pela redução. As demonstrações revelaram a estrutura dos
argumentos que continuarão válidos para extensões mais ricas da
linguagem: o Lema das Formas Canônicas é usado no progresso para
identificar a forma dos valores, e o Lema da Substituição é o passo
técnico essencial da preservação.

A implementação em Haskell do interpretador baseado em ambientes
evidencia a correspondência entre a semântica formal e o código: a
regra (B-Lam) vira a construção de uma closure; a regra (B-App) vira
a sequência de chamadas recursivas com extensão do ambiente. Essa
correspondência é a base metodológica sobre a qual os próximos
capítulos construirão interpretadores para linguagens
progressivamente mais ricas.
